\input{../tex-inputs/latex-korrekturansicht-vorspann}

\section[Theodor Herzl an Arthur Schnitzler, 8. 12. 1893]{L03833 Theodor Herzl an Arthur Schnitzler, 8. 12. 1893}
\nopagebreak\mylabel{L03833v}
\rehead{ }\normalsize\beginnumbering\briefempfaengerindex{Schnitzler, Arthur@\textsc{Schnitzler, Arthur}!zzzHerzl, Theodor@\emph{von Theodor Herzl}!1893-12-082@{8. 12. 1893}|(be}
\toendnotes[C]{\smallbreak\pagebreak[2]}
\correspDesc{Versand  durch Theodor Herzl am 8. 12. 1893 in Paris
\newline{}Erhalt  durch Arthur Schnitzler im Zeitraum [9. 12. 1893
                  – 10. 12. 1893] in Wien}\toendnotes[C]{\smallbreak}
\Standort{CUL, Schnitzler, B 39.}
\physDesc{Brief, 1 Blatt, 1 Seite, 484 Zeichen
\newline{}Handschrift: schwarze Tinte, lateinische Kurrent
\newline{}Ordnung: mit Bleistift von unbekannter Hand nummeriert: »12« }
\buchAbdrucke{\weitereDrucke{1) Theodor Herzl: \emph{Briefe und autobiographische Notizen 1866–1895}. Bearbeitet von Johannes Wachten in Zusammenarbeit mit Chaya Harel, Daisy Tycho und Manfred Winkler. Berlin, Frankfurt am Main, Wien: \emph{Propyläen} 1983, S. 541–542 (Briefe und Tagebücher. Herausgegeben von Alex Bein, Hermann Greive, Moshe Schaerf, Julius H. Schoeps und Johannes Wachten, 1).} \weitereDrucke{2) Hermann Bahr, Arthur Schnitzler: \emph{Briefwechsel, Aufzeichnungen, Dokumente (1891–1931)}. Herausgegeben von Kurt Ifkovits und Martin Anton Müller. (2018) \url{https://schnitzler-bahr.acdh.oeaw.ac.at/L041403.html}.} }\toendnotes[C]{\smallbreak}
\pstart
           {\pb}\textcolor{brown}{\textcolor{gray}{\textbf{\textsc{Nouvelle Presse Libre}}}}\orgindex{Neue Freie Presse@Neue Freie Presse|pw}{}\ledrightnote{\textcolor{brown}{Neue Freie Presse}}\hfill \textcolor{gray}{\textbf{\emph{\textcolor{pink}{Paris}\oindex{8, rue de Monceau@\textbf{8, rue de Monceau}, \emph{Wohngebäude}|pw}{}\ledrightnote{\textcolor{pink}{8, rue de Monceau}}, le}}}{ }8 Dez{ }\textcolor{gray}{\textbf{189}}3\pend
           
\pstart
           \textcolor{gray}{\textbf{D\textsuperscript{r}{ }THÉODORE HERZL}}\pend
           
\pstart{}Lieber Freund!\pend\vspace{0.5em}
\pstart
           Sieben Wochen war ich malariakrank \introOben{}– aus \label{K_L03833-1v}\edtext{\textcolor{pink}{Venedig}\oindex{Venedig@\textbf{Venedig}|pw}{}\ledrightnote{\textcolor{pink}{Venedig}}–\textcolor{pink}{Toulon}\oindex{Toulon@\textbf{Toulon}, \emph{Hauptstadt}|pw}{}\ledrightnote{\textcolor{pink}{Toulon}}}{\lemma{\textnormal{\emph{Venedig–Toulon}}}\Cendnote{\textnormal{\textcolor{blue}{Herzl}\pwindex{Herzl, Theodor 2.\,5.\,1860 Budapest – 3.\,7.\,1904 Edlach@\textsc{Herzl, Theodor} (2.\,5.\,1860 Budapest – 3.\,7.\,1904 Edlach), \emph{Schriftsteller, Journalist}|pwk} war als \textcolor{pink}{Frankreich}\oindex{Frankreich@\textbf{Frankreich}|pwk}korrespondent für die \emph{\textcolor{brown}{Neue Freie Presse}\orgindex{Neue Freie Presse@Neue Freie Presse|pwk}} nach \textcolor{pink}{Toulon}\oindex{Toulon@\textbf{Toulon}, \emph{Hauptstadt}|pwk} gereist und hatte an den Festivitäten zum Empfang der \textcolor{pink}{russischen}\oindex{Russland@\textbf{Russland}|pwk} Flotte dort teilgenommen. In
                     seinem \emph{\textcolor{green}{Bericht}\pwindex{Herzl, Theodor 2.\,5.\,1860 Budapest – 3.\,7.\,1904 Edlach@\textsc{Herzl, Theodor} (2.\,5.\,1860 Budapest – 3.\,7.\,1904 Edlach), \emph{Schriftsteller, Journalist}!Flottenfest in Toulon@\strich\emph{Das Flottenfest in Toulon}|pwk}} darüber erwähnte er auch
                     das \textcolor{pink}{venezianische}\oindex{Venedig@\textbf{Venedig}|pwk} Fest mit »Tausenden von
                     Zuschauern«, das am 14. 10. 1893 an den Kais veranstaltet wurde (\emph{\textcolor{green}{Das Flottenfest in Toulon}\pwindex{Herzl, Theodor 2.\,5.\,1860 Budapest – 3.\,7.\,1904 Edlach@\textsc{Herzl, Theodor} (2.\,5.\,1860 Budapest – 3.\,7.\,1904 Edlach), \emph{Schriftsteller, Journalist}!Flottenfest in Toulon@\strich\emph{Das Flottenfest in Toulon}|pwk}}. In: \emph{\textcolor{green}{Neue Freie Presse}\pwindex{Neue Freie Presse@\emph{Neue Freie Presse}|pwk}}, Nr. 10.471,
                           16. 10. 1893, S. 5).}}}\label{K_L03833-1} mitgebracht –\introOben{}, schwere Fieberwochen. Heute ein
               \label{K_L03833-56v}\edtext{zermatschgerter}{\lemma{\textnormal{\emph{zermatschgerter}}}\Cendnote{\textnormal{umgangssprachlich: zu Brei verdrückter}}}\label{K_L03833-56} Reconvalescent. So fasse ich mich nothgedrungen kurz. Ich gratulire
               zur \label{K_L03833-2v}\edtext{\textcolor{green}{\textcolor{violet}{Märchenaufführung}\eventindex{Volkstheater@\textbf{Volkstheater}!Uraufführung von Das Märchen, 1.12.1893@Uraufführung von Das Märchen, 1.12.1893|pwv}{}\ledrightnote{{$\rightarrow$}\emph{\textcolor{violet}{Uraufführung von Das Märchen, 1.12.1893}}}}\pwindex{Schnitzler, Arthur 15. 5. 1862 Wien – 21. 10. 1931 ebd.@\textsc{Schnitzler, Arthur} (15. 5. 1862 Wien – 21. 10. 1931 ebd.), \emph{Schriftsteller, Mediziner}!Märchen. Schauspiel in drei Aufzügen@\strich\emph{Das Märchen. Schauspiel in drei Aufzügen}|pw}{}\ledrightnote{\textcolor{green}{Das Märchen. Schauspiel in drei Aufzügen}}}{\lemma{\textnormal{\emph{Märchenaufführung}}}\Cendnote{\textnormal{ Die \textcolor{violet}{Uraufführung von \emph{\textcolor{green}{Das Märchen}\pwindex{Schnitzler, Arthur 15. 5. 1862 Wien – 21. 10. 1931 ebd.@\textsc{Schnitzler, Arthur} (15. 5. 1862 Wien – 21. 10. 1931 ebd.), \emph{Schriftsteller, Mediziner}!Märchen. Schauspiel in drei Aufzügen@\strich\emph{Das Märchen. Schauspiel in drei Aufzügen}|pwk}} von \textcolor{blue}{Schnitzler}}\eventindex{Volkstheater@\textbf{Volkstheater}!Uraufführung von Das Märchen, 1.12.1893@Uraufführung von Das Märchen, 1.12.1893|pwk} fand am 1. 12. 1893 am \emph{\textcolor{brown}{Volkstheater}\orgindex{Volkstheater@Volkstheater|pwk}} statt.}}}\label{K_L03833-2}. Die Bühne ist erobert. Sie
               werden sie nicht mehr verlassen. Nur sich treu bleiben! u. ringsherum reden
               lassen.\pend
           
\pstart
           Ich grüsse Sie herzlich {\\[\baselineskip]}Ihr aufrichtig ergebener {\\[\baselineskip]}\spacefill\mbox{Herzl}\pend
           \leftskip=0em{}
\pstart
           \noindent{}Könnten Sie mir \textcolor{blue}{Bahrs}\pwindex{Bahr, Hermann 19.\,7.\,1863 Linz – 15.\,1.\,1934 München@\textsc{Bahr, Hermann} (19.\,7.\,1863 Linz – 15.\,1.\,1934 München), \emph{Schriftsteller, Kritiker}|pw}{}\ledrightnote{\textcolor{blue}{Hermann Bahr}}\label{K_L03833-3v}\edtext{{ }\textcolor{green}{Artikelserie über
                     Jung-Oesterreich}\pwindex{Bahr, Hermann 19.\,7.\,1863 Linz – 15.\,1.\,1934 München@\textsc{Bahr, Hermann} (19.\,7.\,1863 Linz – 15.\,1.\,1934 München), \emph{Schriftsteller, Kritiker}!junge Österreich@\strich\emph{Das junge Österreich}|pwv}{}\ledrightnote{{$\rightarrow$}\emph{\textcolor{green}{Das junge Österreich}}}}{\lemma{\textnormal{\emph{Artikelserie über Jung-Oesterreich}}}\Cendnote{\textnormal{\textcolor{blue}{Hermann Bahr}\pwindex{Bahr, Hermann 19.\,7.\,1863 Linz – 15.\,1.\,1934 München@\textsc{Bahr, Hermann} (19.\,7.\,1863 Linz – 15.\,1.\,1934 München), \emph{Schriftsteller, Kritiker}|pwk}: \emph{\textcolor{green}{Das junge Österreich}\pwindex{Bahr, Hermann 19.\,7.\,1863 Linz – 15.\,1.\,1934 München@\textsc{Bahr, Hermann} (19.\,7.\,1863 Linz – 15.\,1.\,1934 München), \emph{Schriftsteller, Kritiker}!junge Österreich@\strich\emph{Das junge Österreich}|pwk}}. In: \emph{\textcolor{green}{Deutsche Zeitung}\pwindex{Deutsche Zeitung@\emph{Deutsche Zeitung}|pwk}}, Jg. 23, Nr. 7806,
                           20. 9. 1893, Morgen-Ausgabe, S. 1–2; Nr. 7813,
                           27. 9. 1893, Morgen-Ausgabe,
                     S. 1–3; Nr. 7823, 7. 10. 1893, Morgen-Ausgabe,
                        S. 1–3.}}}\label{K_L03833-3} (\textcolor{green}{Deutsche
                  Zeitung}\pwindex{Deutsche Zeitung@\emph{Deutsche Zeitung}|pw}{}\ledrightnote{\textcolor{green}{Deutsche Zeitung}}) verschaffen? Bitte!\pend
           \selectlanguage{ngerman}\endnumbering\briefempfaengerindex{Schnitzler, Arthur@\textsc{Schnitzler, Arthur}!zzzHerzl, Theodor@\emph{von Theodor Herzl}!1893-12-082@{8. 12. 1893}|)be}\mylabel{L03833h}  \newcommand{\dateiname}{L03833}\newcommand{\titel}{Theodor Herzl an Arthur Schnitzler, 8. 12. 1893}\newcommand{\editorInnen}{Selma Jahnke und Martin Anton Müller}\input{../tex-inputs/latex-korrekturansicht-abspann}
